\subsection{Mean-field variational inference}
\label{subsec:mean-field-vi}

\subsubsection{Factorized families and conditional updates}

The ELBO formulation is generic.
The real modeling decision is the choice of variational family \(\mathcal{Q}\).
If the family is too expressive, evaluating or optimizing the ELBO can be almost as hard as the original posterior problem.
If the family is too simple, optimization is cheap but the approximation can be poor.
Mean-field variational inference makes a deliberately simple choice:
it assumes the latent coordinates are independent under the approximation, even if they are strongly dependent under the true posterior.
This is the standard mean-field starting point in variational inference; see \citet[\S2.3]{blei_kucukelbir_mcauliffe_2017_vi_review}.
To see what this means in a hierarchical model, suppose the observations are \(x_{1:n}\), the shared latent block is \(w\), and the datum-specific latent variables are \(z_{1:n}\).
Many posterior distributions of interest take the form
\[
  p_\theta(w,z_{1:n}\mid x_{1:n})
  \propto
  p_\theta(w)\prod_{i=1}^n p_\theta(z_i\mid w)\,p_\theta(x_i\mid z_i,w).
\]
Here the global block \(w\) is shared across the whole data set, while each \(z_i\) is local to one observation.
A basic mean-field approximation replaces this coupled posterior by
\[
  q(w,z_{1:n})
  =
  q(w)\prod_{i=1}^n q_i(z_i),
\]
possibly with a finer factorization inside the global block \(q(w)\).
This suppresses posterior dependence, but it also isolates the global quantities that pool information from all observations from the local quantities that can be updated one datum at a time.
To state the coordinate update in a uniform form, we now collect all latent blocks into \(z=(z_1,\dots,z_m)\) and consider the generic factorized family below.

Write the latent variable as \(z=(z_1,\dots,z_m)\) and consider the factorized family
\begin{equation}
  q(z)
  =
  \prod_{j=1}^m q_j(z_j).
  \label{eq:mean-field-factorization}
\end{equation}
The factorization in \eqref{eq:mean-field-factorization} is strong, but it makes the ELBO amenable to coordinate updates.
To derive one update, we hold all factors except \(q_j\) fixed and optimize only over the remaining factor.
That one-coordinate problem is the basis of coordinate-ascent variational inference (CAVI).
Mean-field VI is often the first approximation one studies because it exposes this coordinatewise logic in the simplest possible setting before one worries about richer families.
The standard derivation and algorithmic viewpoint appear in \citet[\S2.4]{blei_kucukelbir_mcauliffe_2017_vi_review}.

\subsubsection{Coordinate ascent updates}
\label{subsubsec:cavi-updates}

Let
\[
  q_{-j}(z_{-j})
  \coloneqq
  \prod_{k\ne j} q_k(z_k)
\]
denote the product of all factors except the $j$th.
If we hold $q_{-j}$ fixed and optimize only over $q_j$, the ELBO has a simple closed-form maximizer.

\begin{proposition}[Mean-field coordinate update]
\label{prop:mean-field-update}
Fix the factors $q_k$ for $k\ne j$.
Then any maximizer of the ELBO over the remaining factor $q_j$ satisfies
\begin{equation}
  q_j^\star(z_j)
  \propto
  \exp\!\left(
    \E_{q_{-j}}\bigl[\log p_\theta(x,z)\bigr]
  \right),
  \label{eq:mean-field-coordinate-update}
\end{equation}
where the expectation is taken with respect to the current variational law on $z_{-j}$ and the proportionality constant normalizes the right-hand side.
\end{proposition}
\begin{proof}
Starting from \eqref{eq:elbo-definition} and separating the factor $q_j$ from the others,
\[
  \mathcal{L}(q;\theta,x)
  =
  \int q_j(z_j)\,q_{-j}(z_{-j})
  \Bigl[\log p_\theta(x,z) - \log q_j(z_j) - \log q_{-j}(z_{-j})\Bigr]
  \,dz_j\,dz_{-j}.
\]
The terms involving only $q_{-j}$ are constant when we optimize over $q_j$, so
\[
  \mathcal{L}(q;\theta,x)
  =
  \int q_j(z_j)\,m_j(z_j)\,dz_j
  -
  \int q_j(z_j)\log q_j(z_j)\,dz_j
  + C,
\]
where
\[
  m_j(z_j)
  \coloneqq
  \E_{q_{-j}}\bigl[\log p_\theta(x,z)\bigr]
\]
and $C$ does not depend on $q_j$.
Now define
\[
  \widetilde q_j(z_j)
  \propto
  e^{m_j(z_j)}.
\]
Then
\[
  \operatorname{KL}\!\bigl(q_j\,\|\,\widetilde q_j\bigr)
  =
  \int q_j(z_j)\log q_j(z_j)\,dz_j
  -
  \int q_j(z_j)m_j(z_j)\,dz_j
  +
  \log \int e^{m_j(u)}\,du.
\]
Therefore
\[
  \mathcal{L}(q;\theta,x)
  =
  -\operatorname{KL}\!\bigl(q_j\,\|\,\widetilde q_j\bigr)
  + C',
\]
for another constant $C'$ independent of $q_j$.
The ELBO is maximized exactly when this KL divergence is minimized, which occurs at $q_j=\widetilde q_j$.
This gives \eqref{eq:mean-field-coordinate-update}.%
\end{proof}

Equation \eqref{eq:mean-field-coordinate-update} is the core formula behind coordinate-ascent variational inference (CAVI).
One cycles through the factors, replacing each one by the normalized exponential of the expected complete-data log density.
Each update maximizes the ELBO with respect to one coordinate, so the ELBO cannot decrease from one step to the next.
This is the generic CAVI update from \citet[\S2.4, eqs.~(17)--(18)]{blei_kucukelbir_mcauliffe_2017_vi_review}.
In that sense, CAVI is to VI what the E-step/M-step alternation is to EM: each block update solves the current restricted subproblem exactly.
It is also the deterministic analogue of a Gibbs update:
\[
  \text{Gibbs:}\qquad Z_j^{\mathrm{new}} \sim p(z_j\mid z_{-j},x),
  \qquad
  \text{CAVI:}\qquad q_j^{\mathrm{new}}(z_j)\propto \exp\!\left(\E_{q_{-j}}[\log p_\theta(x,z)]\right).
\]
Gibbs sampling draws one coordinate from the exact conditional, while CAVI replaces that random draw by the factor that maximizes the ELBO given the current law on the remaining coordinates.

\begin{algorithm}[H]
  \caption{Coordinate-ascent mean-field VI}
  \label{alg:cavi}
  \begin{algorithmic}[1]
    \Require Joint density \(p_\theta(x,z)\), initial factors $q_1^{(0)},\dots,q_m^{(0)}$, tolerance $\varepsilon$
    \State $t \gets 0$
    \Repeat
      \For{$j=1,2,\dots,m$}
        \State Update
        \[
          q_j^{(t+1)}(z_j)
          \propto
          \exp\!\left(
            \E_{q_{-j}^{(t+1,t)}}[\log p_\theta(x,z)]
          \right),
        \]
        where $q_{-j}^{(t+1,t)}$ uses the newest factors available in the current sweep
      \EndFor
      \State $t \gets t+1$
    \Until{the ELBO increase is smaller than $\varepsilon$}
    \State \Return $q^{(t)}(z)=\prod_{j=1}^m q_j^{(t)}(z_j)$
  \end{algorithmic}
\end{algorithm}

The appeal of CAVI is not only monotone improvement of the ELBO but also a closure property in conjugate models.
Suppose the exact complete conditional of \(z_j\) has exponential-family form
\[
  p_\theta(z_j\mid x,z_{-j})
  \propto
  h_j(z_j)\exp\!\bigl(\eta_j(x,z_{-j})^\top T_j(z_j)\bigr).
\]
Then \eqref{eq:mean-field-coordinate-update} gives
\[
  q_j^\star(z_j)
  \propto
  h_j(z_j)\exp\!\Bigl(\E_{q_{-j}}\bigl[\eta_j(x,z_{-j})\bigr]^\top T_j(z_j)\Bigr).
\]
So the optimal factor stays in the same exponential family; only its natural parameter is replaced by an expectation under the remaining factors.
This is the basic closure principle behind classical mean-field VI:
Gaussian conditionals yield Gaussian factors, categorical conditionals yield categorical factors, and Dirichlet conditionals yield Dirichlet factors.
In conjugate exponential-family models the right-hand side can therefore be evaluated from current sufficient-statistic moments, and each coordinate update reduces to a familiar parameter update rather than a fresh numerical optimization.
This exponential-family closure is the organizing principle in \citet[\S4.1--4.2]{blei_kucukelbir_mcauliffe_2017_vi_review}.
That is what makes mean-field VI computationally attractive in models such as mixture models, hidden Markov models, topic models, and linear-Gaussian factor models.
Once conjugacy disappears, however, those closed-form coordinate updates are no longer available.
That is the point where one leaves the classical CAVI setting and moves to stochastic-gradient and amortized methods.

\subsubsection{Worked examples: mixture assignments and factor analysis}

Two examples are especially useful here.
Gaussian mixtures show how CAVI keeps categorical, Gaussian, and Dirichlet factors inside their familiar conjugate families, while Bayesian factor analysis makes the global/local Gaussian structure visible in a form that leads naturally into amortized encoders and VAEs.

\begin{example}[Gaussian mixtures stay inside familiar families]
\label{ex:mfvi-gmm}
The same closure principle appears in a conjugate Gaussian mixture model.
Let \(z_i\in\{1,\dots,K\}\) be the component indicator for \(x_i\in\R^d\), let \(\pi\) be the mixing weights, and let \(\mu_{1:K}\) be the component means.
Assume the observation covariance is \(I_d\), take the priors
\[
  \pi \sim \operatorname{Dirichlet}(\alpha_1,\dots,\alpha_K),
  \qquad
  \mu_k \sim \mathcal{N}(0,\Lambda^{-1}),
\]
and write the mean-field family as
\[
  q(\pi,\mu_{1:K},z_{1:n})
  =
  q(\pi)\prod_{k=1}^K q_k(\mu_k)\prod_{i=1}^n q_i(z_i).
\]
Then CAVI keeps every factor in its conjugate family:
\[
  q_i(z_i)
  =
  \operatorname{Categorical}(r_{i1},\dots,r_{iK}),
  \qquad
  q_k(\mu_k)
  =
  \mathcal{N}(m_k,\Sigma_k),
  \qquad
  q(\pi)
  =
  \operatorname{Dirichlet}(\widetilde\alpha_1,\dots,\widetilde\alpha_K).
\]
The local responsibility update is
\[
  r_{ik}
  \propto
  \exp\!\left(
    \E_q[\log \pi_k]
    -
    \frac12 \E_q \norm{x_i-\mu_k}^2
  \right),
\]
while the global factors satisfy
\[
  \Sigma_k^{-1}
  =
  \Lambda + \Bigl(\sum_{i=1}^n r_{ik}\Bigr) I_d,
  \qquad
  m_k
  =
  \Sigma_k \sum_{i=1}^n r_{ik}x_i,
\]
and
\[
  \widetilde\alpha_k
  =
  \alpha_k + \sum_{i=1}^n r_{ik}.
\]
This is the packet's basic message in a standard conjugate model:
the local \(z_i\)-updates are soft assignments, and the global updates for \(\mu_k\) and \(\pi\) are just Gaussian and Dirichlet parameter updates driven by the current responsibilities.
\end{example}

\begin{example}[Global loadings and local factors]
\label{ex:mfvi-bfa}
A linear-Gaussian factor model makes the global/local structure explicit.
Let \(x_i\in\R^d\) be observed, let \(z_i\in\R^k\) with \(k<d\) be a latent factor, and let \(W\in\R^{d\times k}\) be a loading matrix.
Write the \(r\)th row of \(W\) as \(w_r^\top\in\R^k\).
With fixed scales \(\sigma^2,\tau^2>0\), consider
\[
  w_r \sim \mathcal{N}(0,\tau^2 I_k),
  \qquad
  z_i \sim \mathcal{N}(0,I_k),
  \qquad
  x_i\mid z_i,W \sim \mathcal{N}(Wz_i,\sigma^2 I_d).
 \]
Then
\[
  p_\theta(W,z_{1:n}\mid x_{1:n})
  \propto
  \prod_{r=1}^d p_\theta(w_r)
  \prod_{i=1}^n p_\theta(z_i)\,p_\theta(x_i\mid z_i,W),
\]
where \(\theta=(\sigma^2,\tau^2)\) collects the fixed hyperparameters.
A mean-field approximation that respects the global/local split is
\[
  q(W,z_{1:n})
  =
  \prod_{r=1}^d q_r(w_r)\prod_{i=1}^n q_i(z_i).
\]
Because both the prior and the likelihood are Gaussian, the closure principle above implies that every optimal factor is Gaussian as well.
In fact one obtains
\[
  q_i(z_i)
  =
  \mathcal{N}(m_i,S),
  \qquad
  q_r(w_r)
  =
  \mathcal{N}(\mu_r,\Sigma),
\]
with shared covariance matrices
\[
  S^{-1}
  =
  I_k + \sigma^{-2}\E_q[W^\top W],
  \qquad
  \Sigma^{-1}
  =
  \tau^{-2}I_k + \sigma^{-2}\sum_{i=1}^n \E_q[z_i z_i^\top],
\]
and means
\[
  m_i
  =
  \sigma^{-2}S\,\E_q[W]^\top x_i,
  \qquad
  \mu_r
  =
  \sigma^{-2}\Sigma\sum_{i=1}^n x_{ir}\,\E_q[z_i].
\]
The required second moments are
\[
  \E_q[W^\top W]
  =
  \sum_{r=1}^d (\Sigma + \mu_r\mu_r^\top),
  \qquad
  \E_q[z_i z_i^\top]
  =
  S + m_i m_i^\top.
\]
This example makes both kinds of mean-field update visible.
Each local factor \(q_i(z_i)\) uses only the datum \(x_i\) together with the current global moments of \(W\).
Each global factor \(q_r(w_r)\) pools information from all observations through the latent second moments \(\E_q[z_i z_i^\top]\).
If \(W\) were held fixed, the local update would reduce to the exact Gaussian posterior from classical factor analysis; the Bayesian version adds one more layer of uncertainty by learning a posterior distribution over the shared loadings themselves.
\end{example}

\subsubsection{Strengths and limitations of mean-field}
\label{subsubsec:mean-field-strengths-limitations}

Mean-field VI is often remarkably effective when the posterior is roughly unimodal and the most important uncertainty is marginal rather than strongly correlated.
Its main weakness is also clear from \eqref{eq:mean-field-factorization}: it cannot represent posterior dependence.
If two coordinates are coupled under the posterior, the factorized approximation must absorb that dependence indirectly by inflating or shrinking the one-dimensional factors.
This is one reason mean-field approximations often understate uncertainty.

The other major limitation is computational scaling with the data set.
In models with one local latent variable per observation, CAVI introduces one local variational factor per observation as well.
That is manageable for moderate data sets and conjugate models, but it becomes cumbersome when $n$ is very large or when the coordinate updates no longer have closed form.
This is exactly the setting that motivates amortized and stochastic variational inference.
From this point onward the ELBO and the model \(p_\theta(x,z)\) stay in place, but the optimization regime changes: instead of approximating the posterior for a fixed \(\theta\), we allow the model and the inference rule to be learned together from data.
